% Part: computability
% Chapter: recursive-functions
% Section: normal-form

\documentclass[../../../include/open-logic-section]{subfiles}

\begin{document}

\olfileid[tr]{cmp}{rec}{nft}
\olsection{Normal Biçim Teoremi}

\begin{thm}[Kleene Normal Biçim Teoremi]
\ollabel{thm:kleene-nf}
Şu özelliğe sahip ilkel özyinelemeli bir $T(e,x,s)$ bağıntısı ve ilkel
özyinelemeli bir $U(s)$ fonksiyonu vardır: $f$ herhangi bir kısmi
özyinelemeli fonksiyonsa bazı $e$ için, her $x$'te
\[
  f(x)\simeq U(\umin{s}{T(e,x,s)})
\]
olur.
\end{thm}

\begin{explain}
Normal biçim teoreminin ispatı ayrıntılıdır, fakat temel düşünce basittir.
Her kısmi özyinelemeli fonksiyonun sezgisel olarak programını ya da tanımını
kodlayan bir sayı olan \emph{indisi}~$e$ vardır. $f(x)\fdefined$ ise hesap
sistematik olarak kaydedilip bazı $s$ sayısıyla kodlanabilir; $s$'nin
$f$ fonksiyonunun $x$ girdisindeki hesabını kodladığı olgusu da yalnızca
$x$ ile $e$ tanımı kullanılarak ilkel özyinelemeli biçimde denetlenebilir.
Dolayısıyla ``indisi $e$ olan fonksiyonun $x$ girdisi için bir hesabı vardır
ve $s$ bu hesabı kodlar'' anlamındaki $T$ bağıntısı ilkel özyinelemelidir.
Hesabın tam kaydı $s$ verildiğinde $s$'nin ``sonucu'', $f(x)$ değeridir ve bu
değer de $s$'den ilkel özyinelemeli biçimde elde edilebilir.

Normal biçim teoremi, herhangi bir kısmi özyinelemeli fonksiyonun tanımı için
tek bir sınırsız aramanın yeterli olduğunu gösterir. Özünde, $e$ indisli
fonksiyonun $x$ girdisindeki hesabını kodlayan bir sayı buluncaya kadar bütün
sayılar arasında arama yapabiliriz. $e$ sayılarını kısmi özyinelemeli
fonksiyonların ``adları'' olarak kullanabilir ve teoremdeki denklemle
tanımlanan $f$ fonksiyonuna $\cfind{e}$ diyebiliriz. Her kısmi özyinelemeli
fonksiyonun birden fazla indisi olabileceğine, hatta her kısmi özyinelemeli
fonksiyonun sonsuz sayıda indisi bulunduğuna dikkat edin.
\end{explain}
\end{document}

